\documentclass[11pt,a4paper]{article}
\usepackage[T1]{fontenc}
\usepackage[utf8]{inputenc}
\usepackage[margin=27mm,headheight=15pt]{geometry}
\usepackage{amsmath,amsthm,amssymb}
\usepackage{newpxtext,newpxmath}
\usepackage{microtype}
\usepackage[dvipsnames]{xcolor}
\usepackage{booktabs,longtable,array}
\usepackage{enumitem}
\usepackage{tikz}
\usepackage[most]{tcolorbox}
\usepackage{listings}
\usepackage{fancyhdr}
\usepackage{aliascnt}
\usepackage{hyperref}
\usepackage[nameinlink,noabbrev]{cleveref}
\definecolor{InkBlue}{HTML}{173D5D}
\definecolor{PaleBlue}{HTML}{F0F5F9}
\definecolor{Muted}{HTML}{52636F}
\hypersetup{colorlinks=true,linkcolor=InkBlue,citecolor=InkBlue,urlcolor=InkBlue,
  pdftitle={Ordinal Order Maps and Finitary Powersets of Ordinal--Finite Grids},
  pdfauthor={Research manuscript prepared with ChatGPT}}
\pagestyle{fancy}
\fancyhf{}
\fancyhead[L]{\small\color{Muted}Ordinal order maps and finitary powersets}
\fancyhead[R]{\small\color{Muted}Research draft}
\fancyfoot[C]{\thepage}
\renewcommand{\headrulewidth}{0.3pt}
\setlength{\parindent}{1.3em}
\setlength{\parskip}{3pt}
\setcounter{tocdepth}{2}
\emergencystretch=1.8em
\newtheorem{theorem}{Theorem}[section]
\newaliascnt{lemma}{theorem}
\newtheorem{lemma}[lemma]{Lemma}
\aliascntresetthe{lemma}
\newaliascnt{proposition}{theorem}
\newtheorem{proposition}[proposition]{Proposition}
\aliascntresetthe{proposition}
\newaliascnt{corollary}{theorem}
\newtheorem{corollary}[corollary]{Corollary}
\aliascntresetthe{corollary}
\theoremstyle{definition}
\newaliascnt{definition}{theorem}
\newtheorem{definition}[definition]{Definition}
\aliascntresetthe{definition}
\newaliascnt{example}{theorem}
\newtheorem{example}[example]{Example}
\aliascntresetthe{example}
\newaliascnt{problem}{theorem}
\newtheorem{problem}[problem]{Problem}
\aliascntresetthe{problem}
\theoremstyle{remark}
\newaliascnt{remark}{theorem}
\newtheorem{remark}[remark]{Remark}
\aliascntresetthe{remark}
\crefname{lemma}{lemma}{lemmas}
\crefname{proposition}{proposition}{propositions}
\crefname{corollary}{corollary}{corollaries}
\crefname{definition}{definition}{definitions}
\crefname{example}{example}{examples}
\crefname{problem}{problem}{problems}
\crefname{remark}{remark}{remarks}
\newcommand{\Mon}{\operatorname{Mon}}
\newcommand{\Pf}{\mathsf{P}_{\!f}}
\newcommand{\Grid}{\mathcal{G}}
\newcommand{\otype}{\operatorname{o}}
\newcommand{\height}{\operatorname{h}}
\newcommand{\width}{\operatorname{w}}
\newcommand{\rk}{\operatorname{rk}}
\newcommand{\op}{\mathrm{op}}
\newcommand{\Id}{\mathcal{J}}
\newcommand{\Ups}{\mathcal{U}}
\newcommand{\Nrep}[2]{\left[#1\right]_{#2}}
\newcommand{\CNF}{\operatorname{CNF}}
\newcommand{\supp}{\operatorname{supp}}
\newcommand{\card}[1]{\lvert #1\rvert}
\newcommand{\set}[1]{\{#1\}}
\newtcolorbox{resultbox}{colback=PaleBlue,colframe=InkBlue,boxrule=0.6pt,
  arc=1pt,left=9pt,right=9pt,top=7pt,bottom=7pt,breakable}
\lstset{basicstyle=\ttfamily\footnotesize,breaklines=true,
  columns=fullflexible,keepspaces=true,frame=single,rulecolor=\color{Muted},
  backgroundcolor=\color{PaleBlue},showstringspaces=false}
\title{\color{InkBlue}\LARGE\bfseries Ordinal Order Maps and Finitary Powersets\\[4pt]
  of Ordinal--Finite Grids\\[10pt]
  \large\normalfont An exact calculation toward an open extension problem}
\author{Research manuscript prepared with ChatGPT}
\date{19 September 2026}

\begin{document}
\maketitle
\thispagestyle{empty}
\begin{abstract}
For a finite poset $Q$ and an arbitrary ordinal $\beta$, let
$\Mon(Q,\beta)$ be the pointwise ordered set of isotone maps from $Q$ to
$\beta$. We establish an explicit finite Cantor-normal-form expression for its
maximal order type. If
$\beta=\omega^{\gamma_0}+\cdots+\omega^{\gamma_{m-1}}$, with the exponents
nonincreasing and repetitions retained, the expression is
\[
 \otype\bigl(\Mon(Q,\beta)\bigr)
 =\bigoplus_{s\in\Mon(Q,m)}
      \omega^{\bigoplus_{q\in Q}\gamma_{s(q)}}.
\]
We also prove that the rank of a map is the natural sum of its coordinates,
and give closed forms for the height. Finitely generated downsets of
$\alpha\times P$, for finite $P$, are identified with
$\Mon(P^{\op},1+\alpha)$; the formulas therefore compute the corresponding
finitary Hoare-powerset invariants. A three-element fork and its dual have
identical finite order-map counts but different transfinite profiles. More
strongly, for each fixed size $n$, a single explicit probe ordinal below
$\omega^\omega$ encodes the entire maximal-order-type profile for all ordinal
targets. Complete proofs, two exact calculation algorithms, and reproducible
finite checks accompany the manuscript.
\end{abstract}

\begin{resultbox}
\textbf{Status and scope.}
The motivating published question asks for extensions of a computable class
of well-quasi-orders. This manuscript settles a concrete family of
calculations in that direction, not the full extension program. The
mathematical arguments are intended as proofs; they have not been formally
verified or independently refereed. A targeted literature search did not
establish priority for these particular formulas. The software checks are
supporting tests, not proofs of the transfinite statements.
\end{resultbox}

\clearpage
\begingroup
\small
\setlength{\parskip}{0pt}
\tableofcontents
\endgroup
\clearpage

\section{The selected problem and the precise outcome}

\subsection{A published open direction}
Abriola, Halfon, Lopez, Schmitz, Schnoebelen, and Vialard study ordinal
invariants of finitary powersets and give an exact calculus for a selected
family of well-quasi-orders. Their concluding section asks how to enlarge
that family while preserving computability of the invariants
\cite[Section~6, ``Elementary wqos'']{AHLLSSV}.
The question is a research direction rather than a single numerical
conjecture. We choose the following explicit subproblem within that direction.

\begin{problem}[Ordinal--finite grid calculation]\label{prob:grid}
Given an ordinal $\alpha$ and a finite poset $P$, compute the maximal order
type and well-founded height of finite subsets of $\alpha\times P$ under
Hoare domination, identifying mutually dominating subsets.
\end{problem}

The product in this problem is the \emph{Cartesian product order}:
\[
 (\xi,p)\leq(\eta,q)
 \quad\Longleftrightarrow\quad \xi\leq\eta\ \text{and}\ p\leq_Pq.
\]
No lexicographic order is intended. The finite factor is arbitrary: it need
not be a chain, antichain, tree, or series--parallel poset.

The reason this family is approachable is structural. A finitely generated
downset has, above each point of $P$, either an empty ordinal fiber or an
initial segment with a largest point. Thus an infinite set-theoretic problem
reduces to finitely many ordinal coordinates linked by monotonicity
constraints. The remaining difficulty is not cardinality: a linear extension
can place different infinite coordinate patterns in different ordinal
positions, and ordinal addition can erase earlier terms.

\subsection{What is proved here}
The main formula is \cref{thm:block}, whose proof works for every ordinal
$\beta$. Its key ingredient is the pure-monomial calculation
\[
 \otype\bigl(\Mon(Q,\omega^\delta)\bigr)
 =\omega^{\Nrep{\delta}{|Q|}},
\]
where $\Nrep{\delta}{n}$ is the natural sum of $n$ copies of $\delta$.
It is essential that this exponent is generally \emph{not} the ordinary
product $\delta\cdot n$.

The rank formula in \cref{thm:rank} is independent of the maximal-order-type
argument. The exact downset representation in \cref{thm:reduction} transfers
both results to \cref{prob:grid}. The fork example in \cref{sec:fork} shows
why finite order-map counts alone lose relevant information. Finally,
\cref{thm:probe} identifies a finite ordered-fiber signature and proves that
one carefully chosen ordinal target recovers it completely.

This is a complete calculation for the stated family, but it does not
compute arbitrary iterated finitary powersets or the ordinal antichain width
of all these map spaces. In particular, it does not turn the larger
published extension question into a solved problem merely by adding a new
symbol to a grammar.

\section{Conventions and standard ingredients}

\subsection{Orders, ranks, and linear extensions}
A poset is a \emph{well partial order} (wpo) if every infinite sequence has
an increasing pair. All orders considered below are sets, and the
mathematical arguments are in ZFC. An ordinal is identified with its
well-order of predecessors.

For a well-founded poset $X$, define the point rank and height by
\begin{equation}\label{eq:rankdef}
 r_X(x)=\sup_{y<x}\bigl(r_X(y)+1\bigr),
 \qquad
 \height(X)=\sup_{x\in X}\bigl(r_X(x)+1\bigr).
\end{equation}
Thus $\height(\varnothing)=0$, while a singleton has height $1$.
For a wpo $X$, its maximal order type $\otype(X)$ is the maximum of the
order types of its linear extensions. This maximum exists by the
maximal-linearization theorem of de Jongh and Parikh
\cite{deJonghParikh}; see also \cite[Theorem~3.4]{DSS}.

We use the following standard identities, with $\sqcup$ denoting disjoint
union with no cross-comparabilities:
\begin{equation}\label{eq:standard}
 \otype(X\sqcup Y)=\otype(X)\oplus\otype(Y),
 \qquad
 \otype(X\times Y)=\otype(X)\otimes\otype(Y).
\end{equation}
The product identity is the Cartesian-product theorem, not a theorem about
lexicographic products \cite[Theorem~3.5]{deJonghParikh};
\cite[Lemmas~4.2 and~4.5]{DSS}. Induced subposets cannot have larger maximal
order type. Adding comparisons on a fixed underlying set also cannot
increase maximal order type: every linear extension of the augmented order
is already a linear extension of the original one.

For the small examples where width is used, $\width(X)$ denotes the rank
of the tree of finite sequences of pairwise incomparable elements. If
antichains have maximum finite size $k$, this ordinal width is exactly $k$.
We do not confuse ordinal width with a supremum of antichain cardinalities
when the latter is unbounded.

\subsection{Ordinary and natural arithmetic}
Ordinary ordinal sum and product are written $+$ and $\cdot$.
Hessenberg natural sum and product are written $\oplus$ and $\otimes$.
If ordinals are written over a common decreasing sequence of exponents,
\[
 a=\sum_i\omega^{\delta_i}a_i,
 \qquad b=\sum_i\omega^{\delta_i}b_i,
\]
then $a\oplus b=\sum_i\omega^{\delta_i}(a_i+b_i)$. Natural product is
polynomial multiplication with natural addition of exponents:
\[
 \left(\sum_i\omega^{\delta_i}a_i\right)
 \otimes
 \left(\sum_j\omega^{\varepsilon_j}b_j\right)
 =\bigoplus_{i,j}\omega^{\delta_i\oplus\varepsilon_j}(a_i b_j).
\]
Both operations are commutative and associative; natural sum is strictly
increasing in either argument. All natural sums used in our formulas are
finite.

To make finite repetition unambiguous, define
\begin{equation}\label{eq:repeat}
 \Nrep{\delta}{n}
 :=\underbrace{\delta\oplus\cdots\oplus\delta}_{n\text{ copies}}
 =\delta\otimes n,
 \qquad \Nrep{\delta}{0}=0.
\end{equation}
Thus $\Nrep{\omega+1}{2}=\omega\cdot2+2$, whereas
$(\omega+1)\cdot2=\omega\cdot2+1$.
For finite exponents, natural and ordinary finite sums coincide.

Every nonzero ordinal has a finite Cantor normal form. We often expand
its finite coefficients into repetitions:
\begin{equation}\label{eq:beta}
 \beta=\omega^{\gamma_0}+\cdots+\omega^{\gamma_{m-1}},
 \qquad \gamma_0\geq\cdots\geq\gamma_{m-1}.
\end{equation}
For $\beta=0$, take $m=0$. The exponents may be arbitrary ordinals; no
assumption such as $\beta<\varepsilon_0$ is made in the mathematics.

\subsection{Map posets and finite powersets}
For a finite poset $Q$, let
\[
 \Mon(Q,\beta)
 =\{f:Q\to\beta: p\leq_Qq\Rightarrow f(p)\leq f(q)\},
\]
ordered pointwise. It is an induced subposet of $\beta^{|Q|}$ and hence a
wpo. If $Q$ is empty, there is exactly one map, including when $\beta=0$.
If $Q$ is nonempty and $\beta=0$, there are no maps.

For a poset $A$, let $\Pf(A)$ include the empty finite subset and equip it
with the Hoare relation
\begin{equation}\label{eq:Hoare}
 S\leq_H T
 \quad\Longleftrightarrow\quad
 \forall s\in S\ \exists t\in T\ (s\leq_A t).
\end{equation}
This is a quasi-order. Write $S\equiv_H T$ if both directions hold, and
put
\[
 \Grid(\alpha,P):=\Pf(\alpha\times P)/{\equiv_H}.
\]
All powerset invariants in this manuscript refer to this quotient, or
equivalently to the quasi-order with equivalent points identified.

\section{Finitely generated downsets are ordinal order maps}

\begin{lemma}\label{lem:downsetquotient}
For finite subsets $S,T$ of any poset $A$,
\[
 S\leq_HT\quad\Longleftrightarrow\quad\downarrow S\subseteq\downarrow T.
\]
Consequently $\Pf(A)/{\equiv_H}$ is isomorphic to the inclusion order on
finitely generated downsets of $A$.
\end{lemma}
\begin{proof}
If $S\leq_HT$ and $x\leq s\in S$, choose $t\in T$ with $s\leq t$.
Then $x\leq t$, proving inclusion. Conversely, if the downset inclusion
holds, every $s\in S$ lies below a generator in $T$. Equality of generated
downsets is therefore exactly mutual Hoare domination.
\end{proof}

\begin{theorem}[Endpoint representation]\label{thm:reduction}
For every ordinal $\alpha$ and every finite poset $P$,
\begin{equation}\label{eq:reduction}
 \boxed{\Grid(\alpha,P)\cong\Mon(P^{\op},1+\alpha).}
\end{equation}
The $1+$ on the right is ordinary ordinal addition on the left of
$\alpha$.
\end{theorem}
\begin{proof}
Let $D=\downarrow S\subseteq\alpha\times P$, where $S$ is finite. For
$p\in P$, consider
\[
 D_p:=\{\xi<\alpha:(\xi,p)\in D\}.
\]
The fiber $D_p$ is the union of the finitely many principal initial
segments generated by coordinates $\eta$ for which $(\eta,q)\in S$ and
$p\leq_Pq$. If there are no such generators, it is empty. Otherwise the
finite set of these coordinates has a maximum $a_p$, and
$D_p=[0,a_p]$.

Introduce the ordered alphabet $B_\alpha=\{\bot\}+\alpha$, with a new
least element $\bot$. It has order type $1+\alpha$. Encode $D$ by the map
$e_D:P\to B_\alpha$ that sends an empty fiber to $\bot$ and a nonempty
fiber to its largest point $a_p$. If $p\leq_Pq$, downward closure gives
$D_q\subseteq D_p$, so $e_D(q)\leq e_D(p)$. Thus $e_D$ is antitone on
$P$, or isotone on $P^{\op}$.

Conversely, given an antitone map $e:P\to B_\alpha$, let
\[
 D(e)=\{(\xi,p): e(p)\ne\bot\text{ and }\xi\leq e(p)\}.
\]
This is a downset: if $(\xi,p)$ belongs to it and
$(\eta,q)\leq(\xi,p)$, then $q\leq_Pp$ implies
$e(q)\geq e(p)\geq\xi\geq\eta$. Moreover it is generated by the finite
set of pairs $(e(p),p)$ for which $e(p)\ne\bot$.

The two constructions are inverse. Fiberwise inclusion is precisely
pointwise comparison in $B_\alpha$, proving the order isomorphism.
\end{proof}

For finite $\alpha=k$, the target alphabet has $k+1$ elements. For infinite
$\alpha$, one has $1+\alpha=\alpha$. The distinction is not cosmetic:
for example $\Grid(1,P)$ is the finite ideal lattice of $P$, whereas
$\Mon(P^{\op},1)$ is only a singleton. Also, maximal elements $\omega$
in a fiber of $(\omega+1)\times P$ must not be deleted by replacing
$\omega+1$ with $\omega$.

The representation already explains the dual appearing throughout the
paper. A set of positions with a \emph{smaller} endpoint is upward closed
in $P$. Replacing these filters by ideals of $P$ would interchange the
fork examples later and give the wrong answer.

\section{The pure-monomial calculation}

\subsection{A lexicographic construction for one exponent term}
\begin{lemma}\label{lem:lex}
Let $Q$ have $n$ elements, and let $\delta>0$. There is a linear extension
of $\Mon(Q,\omega^\delta)$ of order type
$(\omega^\delta)^n=\omega^{\delta\cdot n}$.
\end{lemma}
\begin{proof}
Choose a topological listing $q_1,\ldots,q_n$ of $Q$: if $q_i<_Qq_j$,
then $i<j$. Compare maps lexicographically in the coordinate sequence
$f(q_1),\ldots,f(q_n)$, with the earliest differing coordinate decisive.
A pointwise strict comparison is respected, so this is a linear extension.

After any admissible prefix, the next coordinate must be at least the
maximum value at its already assigned strict predecessors. There are no
constraints from unassigned predecessors, because the listing is
topological. All values in the resulting tail $[b,\omega^\delta)$ are
allowed, and every such prefix has a completion. The tail has order type
$\omega^\delta$, since this ordinal is additively indecomposable and
$\delta>0$.

Induct backward on the number of remaining coordinates. With no remaining
coordinates there is one completion. If every completion fiber after the
next choice has type $(\omega^\delta)^k$, their lexicographic concatenation
over the tail of possible next values has type
$(\omega^\delta)^k\cdot\omega^\delta=(\omega^\delta)^{k+1}$.
At the empty prefix the order type is $(\omega^\delta)^n$.
\end{proof}

The construction is not always maximal. For two coordinates and
$\delta=\omega+1$, it gives $\omega^{\omega\cdot2+1}$, whereas the
maximal type will be $\omega^{\omega\cdot2+2}$. The missing term comes
from the lower-order part of the exponent. This is why merely sorting
coordinates is insufficient.

\subsection{Splitting the exponent rather than taking a limit}
\begin{theorem}[Monomial theorem]\label{thm:monomial}
For every finite poset $Q$ of size $n$ and every ordinal $\delta$,
\begin{equation}\label{eq:monomial}
 \boxed{\otype\bigl(\Mon(Q,\omega^\delta)\bigr)
       =\omega^{\Nrep{\delta}{n}}.}
\end{equation}
\end{theorem}
\begin{proof}
Since the map space is an induced subposet of $(\omega^\delta)^n$ with
its Cartesian product order, \eqref{eq:standard} gives
\[
 \otype\bigl(\Mon(Q,\omega^\delta)\bigr)
 \leq\underbrace{\omega^\delta\otimes\cdots\otimes\omega^\delta}_{n}
 =\omega^{\Nrep{\delta}{n}}.
\]
It remains to construct a linear extension attaining the bound. If $n=0$
or $\delta=0$, the map space is a singleton and the assertion follows.
Assume henceforth that $n>0$ and $\delta>0$.

We induct on the number of \emph{distinct terms in the Cantor normal form
of $\delta$}, proving the assertion simultaneously for all finite posets.
If $\delta=\omega^\rho c$ has only one term, where $c$ is a positive
integer, then
\[
 \delta\cdot n=\omega^\rho(cn)=\Nrep{\delta}{n}.
\]
The extension in \cref{lem:lex} therefore attains the upper bound.

Suppose $\delta$ has at least two terms. Write
\[
 \delta=\eta+\lambda,
 \qquad \lambda=\omega^\rho c,
\]
where $\lambda$ is its last term and every exponent occurring in the
Cantor normal form of $\eta$ is strictly larger than $\rho$. Set
\[
 \mu=\omega^\eta,\qquad \nu=\omega^\lambda.
\]
Then $\omega^\delta=\mu\cdot\nu$. Every coordinate value has a unique
representation $\mu u+v$ with $u<\nu$ and $v<\mu$.

If $f\in\Mon(Q,\mu\nu)$, its quotient coordinates $u_f(q)$ form an
isotone map $u_f:Q\to\nu$. Fix one such quotient map $u$. Define a finite
poset $Q_u$ on the same underlying set as $Q$ by retaining exactly the
relations
\[
 p\leq_{Q_u}q
 \quad\Longleftrightarrow\quad
 p\leq_Qq\ \text{and}\ u(p)=u(q).
\]
This is the disjoint union of the induced subposets on the equal-value
fibers of $u$. It still has $n$ elements.

For coordinates with distinct quotient values, the order of the ordinal
blocks already decides the comparison, regardless of the remainders.
For equal quotient values, the original comparison is equivalent to a
comparison of remainders. Consequently the fiber over $u$ is isomorphic to
$\Mon(Q_u,\mu)$. The induction hypothesis applies to $\eta$, which has
one fewer Cantor term, and gives the same maximal type for every fiber:
\[
 A=\omega^{\Nrep{\eta}{n}}.
\]
Choose an attaining linear extension in each fiber. There are only
finitely many possible relations $Q_u$, so such choices can also be made
by first choosing an extension for each possible finite relation.

By the one-term case, the outer map space $\Mon(Q,\nu)$ has a linear
extension of type
\[
 B=\omega^{\Nrep{\lambda}{n}}.
\]
Concatenate the chosen fiber extensions in that outer order. If $f\leq g$
in the original map poset, then $u_f\leq u_g$ pointwise. Distinct quotient
fibers are therefore placed in a compatible order, and comparisons inside
a fiber are respected by construction. Hence this concatenation is a
linear extension of the full map poset, with type
\[
 A\cdot B
 =\omega^{\Nrep{\eta}{n}+\Nrep{\lambda}{n}}
 =\omega^{\Nrep{\delta}{n}}.
\]
The last equality holds because the Cantor supports of $\eta$ and
$\lambda$ occupy disjoint, correctly ordered exponent ranges. This matches
the upper bound and completes the induction.
\end{proof}

\begin{remark}[Why this handles arbitrary ordinals]
The proof uses the finite length of the Cantor normal form of the exponent,
not a hereditary notation system and not a cofinal sequence. It therefore
also applies to uncountable exponents and to fixed points such as
$\varepsilon_0$. There is no appeal to an unproved continuity law for
natural sum. Indeed,
$\sup_{k<\omega}(k\oplus k)=\omega$ but
$\omega\oplus\omega=\omega\cdot2$.
\end{remark}

\begin{example}[A mixed exponent]\label{ex:mixed}
For any two-element poset $Q$ and
$\beta=\omega^{\omega+1}=\omega^\omega\cdot\omega$,
\[
 \otype(\Mon(Q,\beta))=\omega^{\omega\cdot2+2}.
\]
The quotient maps into $\omega$ can be linearized in type $\omega^2$.
Every remainder fiber, regardless of the retained finite relation, has
maximal type $\omega^{\omega\cdot2}$. Their concatenation has type
$\omega^{\omega\cdot2}\cdot\omega^2$, giving the asserted exponent.
\end{example}

\section{The finite Cantor-block formula}

\begin{theorem}[Order-map formula]\label{thm:block}
Let $Q$ be finite and write $\beta$ as in \eqref{eq:beta}. For an isotone
block assignment $s:Q\to m$, define
\[
 E(s)=\bigoplus_{q\in Q}\gamma_{s(q)}.
\]
Then
\begin{equation}\label{eq:main}
 \boxed{\otype\bigl(\Mon(Q,\beta)\bigr)
       =\bigoplus_{s\in\Mon(Q,m)}\omega^{E(s)}.}
\end{equation}
The usual empty-map and empty-sum conventions make this valid also when
$Q=\varnothing$ or $\beta=0$.
\end{theorem}
\begin{proof}
Let $b_j=\sum_{i<j}\omega^{\gamma_i}$. Every value in the $j$th ordinal
block has a unique representation $b_j+v$ with $v<\omega^{\gamma_j}$.
An isotone map $f:Q\to\beta$ therefore determines an isotone map $s_f:Q\to m$
recording its block indices.

For a fixed block assignment $s$, the fiber $X_s$ is isomorphic to
\begin{equation}\label{eq:blockfiber}
 X_s\cong
 \prod_{j<m}\Mon\bigl(Q|_{s^{-1}(j)},\omega^{\gamma_j}\bigr).
\end{equation}
Indeed, an original comparison within one block constrains the remainder
coordinates, and a comparison between distinct blocks is already decided.
The monomial theorem and the Cartesian-product identity give
\[
 \otype(X_s)
 =\bigotimes_{j<m}\omega^{\Nrep{\gamma_j}{|s^{-1}(j)|}}
 =\omega^{E(s)}.
\]

If all comparisons between distinct fibers are deleted, one obtains their
finite disjoint union. Adding the actual cross-fiber comparisons cannot
increase maximal order type, so
\[
 \otype(\Mon(Q,\beta))
 \leq\bigoplus_{s\in\Mon(Q,m)}\omega^{E(s)}.
\]
An upper bound alone is not sufficient: it remains to show that the fibers
can be ordered without ordinal absorption.

If $s\leq t$ pointwise, then $\gamma_{s(q)}\geq\gamma_{t(q)}$ for every
$q$, because the block exponents are nonincreasing. Hence
\begin{equation}\label{eq:reverseweight}
 s\leq t\quad\Longrightarrow\quad E(s)\geq E(t).
\end{equation}
List the finitely many block assignments in decreasing order of $E$.
Within a group with the same $E$, choose a linear extension of their
induced pointwise order. This listing respects every pointwise comparison
of block assignments by \eqref{eq:reverseweight}.

Now concatenate maximal linear extensions of $X_s$ in that listing.
A comparison $f\leq g$ across different fibers implies $s_f\leq s_g$,
so the concatenation respects it. Its order type is the ordinary sum of
the monomials $\omega^{E(s)}$ in nonincreasing exponent order. Such an
ordinary sum equals their natural sum. Thus the upper bound is attained.
\end{proof}

\begin{corollary}[Exact grid formula]\label{cor:gridformula}
Write $1+\alpha=\sum_{j<m}\omega^{\gamma_j}$ with repeated nonincreasing
exponents. Then
\begin{equation}\label{eq:gridformula}
 \otype(\Grid(\alpha,P))
 =\bigoplus_{s\in\Mon(P^{\op},m)}
       \omega^{\bigoplus_{p\in P}\gamma_{s(p)}}.
\end{equation}
\end{corollary}
\begin{proof}
Apply \cref{thm:reduction,thm:block}.
\end{proof}

\subsection{Extremal checks and the leading term}
For a chain $C_n$ and an antichain $A_n$ of size $n$,
\begin{align}
 \otype(\Mon(C_n,\beta))
 &=\bigoplus_{a_0+\cdots+a_{m-1}=n}
      \omega^{\bigoplus_{j<m}\Nrep{\gamma_j}{a_j}},\label{eq:chain}\\
 \otype(\Mon(A_n,\beta))&=\beta^{\otimes n}.\label{eq:anti}
\end{align}
The first formula has one assignment per histogram, because a weakly
increasing sequence is determined by the multiplicities of its values.
The second is the Cartesian-product formula. Imposing a total order
extending $Q$ gives the lower bound, and deleting all comparisons gives
the upper bound
\begin{equation}\label{eq:bounds}
 \otype(\Mon(C_n,\beta))
 \leq\otype(\Mon(Q,\beta))\leq\beta^{\otimes n}.
\end{equation}

Let $\Omega_Q(c)=|\Mon(Q,c)|$ for a finite chain $c$. If the leading term
of a nonzero $\beta$ is $\omega^\gamma c$ and $n>0$, then the leading term
of \eqref{eq:main} is
\begin{equation}\label{eq:leading}
 \omega^{\Nrep{\gamma}{n}}\cdot\Omega_Q(c).
\end{equation}
Exactly the assignments using only the $c$ leading equal-exponent blocks
have the largest exponent. Assigning even one coordinate to a strictly
smaller exponent strictly lowers the natural sum. This also covers finite
$\beta$, where the full result is the finite ordinal $\Omega_Q(\beta)$.

\section{Point ranks and exact height}

The maximal-order-type formula retains detailed information about the
finite poset. Height is strikingly simpler: it depends only on the number
of coordinates and the target ordinal. The essential point is an exact
rank formula, not merely the existence of a cofinal diagonal.

\begin{theorem}[Rank of an isotone map]\label{thm:rank}
For every finite $Q$, ordinal $\beta$, and $f\in\Mon(Q,\beta)$,
\begin{equation}\label{eq:rankformula}
 \boxed{r(f)=\bigoplus_{q\in Q}f(q).}
\end{equation}
\end{theorem}
\begin{proof}
Write $R(f)=\bigoplus_q f(q)$ and induct on this ordinal. If $g<f$
pointwise, strict monotonicity of natural sum gives $R(g)<R(f)$.
Therefore all such $g$ are covered by the induction hypothesis and
\[
 r(f)=\sup_{g<f}\bigl(R(g)+1\bigr)\leq R(f).
\]
The zero map has rank $0$. Suppose that at least one coordinate is nonzero.
Choose the smallest trailing Cantor exponent $\delta$ among all nonzero
values $f(q)$. Choose one value $a$ with trailing exponent $\delta$, and
choose $q$ minimal in its equal-value fiber $f^{-1}(a)$. Every strict
predecessor of $q$ then has value strictly below $a$: isotonicity gives
at most $a$, and equality is excluded by minimality.

Because there are finitely many predecessors, their values have a common
bound $b<a$; when there are none, take $b=0$. Replacing only $f(q)=a$ by
any $a'$ with $b\leq a'<a$ preserves isotonicity. Comparisons from
predecessors remain valid, and comparisons to successors remain valid
because their values were at least $a$.

If $\delta=0$, the ordinal $a$ is a successor. Replace it by its
predecessor. This is at least $b$, and subtracts exactly one from the
finite coefficient of the total natural sum. Thus there is $g<f$ with
$R(g)+1=R(f)$, proving the required lower bound.

Suppose $\delta>0$. Write
\[
 a=A+\omega^\delta c,\qquad c>0,
\]
where $A$ has only exponents larger than $\delta$. For $\eta<\omega^\delta$
put
\[
 a_\eta'=A+\omega^\delta(c-1)+\eta.
\]
These values are cofinal in $a$, so all sufficiently large choices satisfy
$a_\eta'\geq b$. Since no other nonzero coordinate has a term below
$\delta$, the total sum has the form
\[
 R(f)=B+\omega^\delta C,\qquad C>0,
\]
and the modified map $g_\eta$ satisfies
\[
 R(g_\eta)=B+\omega^\delta(C-1)+\eta.
\]
Consequently
$\sup_\eta(R(g_\eta)+1)=R(f)$, even after restricting to the sufficiently
large admissible $\eta$. By induction $r(g_\eta)=R(g_\eta)$, so
$r(f)\geq R(f)$. This proves equality in all cases.
\end{proof}

\begin{theorem}[Height formula]\label{thm:height}
Let $n=|Q|>0$. Then
\begin{equation}\label{eq:heightsup}
 \boxed{\height(\Mon(Q,\beta))
       =\sup_{\xi<\beta}\bigl(\Nrep{\xi}{n}+1\bigr).}
\end{equation}
In particular:
\begin{align}
 \height(\Mon(Q,0))&=0,\label{eq:hzero}\\
 \height(\Mon(Q,\zeta+1))&=\Nrep{\zeta}{n}+1,\label{eq:hsucc}\\
 \height(\Mon(Q,\rho+\omega^\delta))
   &=\Nrep{\rho}{n}+\omega^\delta
       \quad(\delta>0),\label{eq:hlimit}
\end{align}
where in \eqref{eq:hlimit} $\omega^\delta$ is one copy of the last Cantor
monomial of the nonzero limit target. For $n=0$, the height is always $1$.
\end{theorem}
\begin{proof}
Every map $f$ has a largest coordinate value $\xi<\beta$, and then
$R(f)\leq\Nrep{\xi}{n}$. Conversely, for every $\xi<\beta$ the constant
map with value $\xi$ is isotone and has rank $\Nrep{\xi}{n}$. Equation
\eqref{eq:heightsup} follows from \cref{thm:rank}. For a successor target,
the supremum is attained at $\zeta$, proving \eqref{eq:hsucc}.

For a nonzero limit, write its last term as $\omega^\delta c$ with
$\delta>0$, and put the first $c-1$ copies into $\rho$. Values
$\rho+\eta$, $\eta<\omega^\delta$, are cofinal in the target. All
exponents in $\rho$ are at least $\delta$, whereas those in $\eta$ are
strictly smaller. Hence
\[
 \Nrep{\rho+\eta}{n}=\Nrep{\rho}{n}+\Nrep{\eta}{n}.
\]
For $0<n<\omega$, each $\Nrep{\eta}{n}<\omega^\delta$, and these
ordinals are cofinal in $\omega^\delta$ because
$\Nrep{\eta}{n}\geq\eta$. Taking the supremum after adding $1$ proves
\eqref{eq:hlimit}.
\end{proof}

For example, with $n=3$ and any three-element poset $Q$,
\[
\begin{array}{c|c}
\beta&\height(\Mon(Q,\beta))\\\hline
\omega&\omega\\
\omega+1&\omega\cdot3+1\\
\omega\cdot2&\omega\cdot4\\
\omega^2+\omega\cdot2&\omega^2\cdot3+\omega\cdot4
\end{array}
\]
The apparently unusual coefficient $4$ in the third row is correct:
$\xi=\omega+k$ gives rank $\omega\cdot3+3k$, whose successors are cofinal
in $\omega\cdot4$.

Combining \cref{thm:reduction,thm:height} computes
$\height(\Grid(\alpha,P))$ by substituting $\beta=1+\alpha$. For finite
$\alpha=k$ and $|P|=n>0$, the answer is $nk+1$, agreeing with the height
of the ideal lattice of a finite $nk$-element grid.

\section{Specializations and a small separation example}

\subsection{Repeated equal monomials}
If $\beta=\omega^\delta\cdot c$ for a positive integer $c$, every block
assignment has exponent $\Nrep{\delta}{n}$. Thus
\begin{equation}\label{eq:equal}
 \otype(\Mon(Q,\omega^\delta\cdot c))
 =\omega^{\Nrep{\delta}{n}}\cdot\Omega_Q(c).
\end{equation}
The coefficient is an ordinary finite order-map count. A total order
extending $Q$ and the unconstrained antichain give
\[
 \binom{c+n-1}{n}\leq\Omega_Q(c)\leq c^n.
\]
The lower count enumerates weakly increasing sequences of length $n$ from
$c$ values; the upper count enumerates all functions.

For $\delta>0$, the same expression computes
$\otype(\Grid(\omega^\delta\cdot c,P))$ with $Q=P^{\op}$.
The counts $\Omega_P(c)$ and $\Omega_{P^{\op}}(c)$ agree: compose a map
with the reversal $j\mapsto c-1-j$. When $\delta=0$ and the grid ordinal
is finite, one must instead use the target $c+1$.

\subsection{Filters at the first successor of a limit}
For a finite poset $P$ of size $n$, define its filter-size polynomial by
\begin{equation}\label{eq:filterpoly}
 U_P(t)=\sum_{F\in\Ups(P)}t^{|F|}
       =\sum_{k=0}^n u_k t^k,
\end{equation}
where $\Ups(P)$ includes both the empty filter and all of $P$.

\begin{corollary}\label{cor:omegaone}
For $|P|=n$,
\begin{equation}\label{eq:omegaone}
 \boxed{\otype(\Grid(\omega+1,P))
 =\omega^n u_n+\omega^{n-1}u_{n-1}+\cdots+\omega u_1+u_0.}
\end{equation}
For $n>0$, its height is $\omega\cdot n+1$.
\end{corollary}
\begin{proof}
The two blocks of $\omega+1$ have exponents $1$ and $0$. An isotone
assignment from $P^{\op}$ to $2$ is determined by its zero-fiber, which
is an ideal in $P^{\op}$, hence a filter $F$ in $P$. The exponent weight
is exactly $|F|$. Group the main formula by this size. The height follows
from \eqref{eq:hsucc}.
\end{proof}

For a chain, $u_k=1$ for every $k$. For an antichain,
$u_k=\binom nk$. Thus the formulas interpolate between
$\omega^n+\cdots+\omega+1$ and the natural power
$(\omega+1)^{\otimes n}$.

\subsection{The fork and its dual}\label{sec:fork}
Let $V$ have elements $a,b,c$ with $a<b$ and $a<c$, and no comparison
between $b$ and $c$. Let $\Lambda=V^{\op}$.

\begin{center}
\begin{tikzpicture}[scale=0.95,every node/.style={font=\small}]
\node (a) at (0,0) {$a$}; \node (b) at (-.7,1) {$b$};
\node (c) at (.7,1) {$c$};
\draw (a)--(b) (a)--(c);
\node at (0,-.55) {$V$};
\node (aa) at (4,1) {$a$}; \node (bb) at (3.3,0) {$b$};
\node (cc) at (4.7,0) {$c$};
\draw (aa)--(bb) (aa)--(cc);
\node at (4,-.55) {$\Lambda=V^{\op}$};
\end{tikzpicture}
\end{center}

The filters of $V$ are
$\varnothing,\{b\},\{c\},\{b,c\},\{a,b,c\}$. Those of $\Lambda$ are
$\varnothing,\{a\},\{a,b\},\{a,c\},\{a,b,c\}$. Hence
\[
 U_V(t)=1+2t+t^2+t^3,
 \qquad U_\Lambda(t)=1+t+2t^2+t^3,
\]
and \cref{cor:omegaone} yields
\begin{align}
 \otype(\Grid(\omega+1,V))
 &=\omega^3+\omega^2+\omega\cdot2+1,\label{eq:forkV}\\
 \otype(\Grid(\omega+1,\Lambda))
 &=\omega^3+\omega^2\cdot2+\omega+1.\label{eq:forkL}
\end{align}
Their heights are both $\omega\cdot3+1$.

\begin{proposition}\label{prop:separation}
The two underlying grids $(\omega+1)\times V$ and
$(\omega+1)\times\Lambda$ have the same triple of standard invariants,
\[
 (\otype,\height,\width)=(\omega\cdot3+3,\ \omega+2,\ 3),
\]
but their finitary Hoare powersets have different maximal order types.
Moreover $\Omega_V(k)=\Omega_\Lambda(k)$ for every finite $k$.
\end{proposition}
\begin{proof}
Any three-element finite poset has maximal order type $3$. The Cartesian
product formula gives
$(\omega+1)\otimes3=\omega\cdot3+3$ for both grids. Each finite fork has
height $2$. The rank in a Cartesian product is the natural sum of its
point ranks \cite[Lemma~4.6]{DSS}. The largest grid point rank is therefore
$\omega\oplus1=\omega+1$, and the height is $\omega+2$.

Each grid is the union of three ordinal chains, one over each point of
the finite factor, so an antichain has at most three points. More
generally, if $p_1,\ldots,p_n$ is a topological listing of a finite $P$,
the points $(n-i,p_i)$, $1\leq i\leq n$, are pairwise incomparable in
$\omega\times P$. A comparison in the finite factor forces the ordinal
coordinates to run in the opposite direction. This proves that both
widths here are exactly $3$.

The powerset types differ by \eqref{eq:forkV}--\eqref{eq:forkL}. Finally,
value reversal gives a bijection between isotone maps from a finite poset
and its dual to a finite chain. In this example the common count is
\[
 \Omega_V(k)=\sum_{a=0}^{k-1}(k-a)^2
 =\sum_{j=1}^k j^2.
\]
The formula follows by choosing the value at the minimum and then choosing
the two upper values independently.
\end{proof}

Thus equality of all finite order-map counts does not determine the
transfinite maximal-order-type profile. This is not presented as the first
nonfunctionality phenomenon for finitary powersets; such nonfunctionality
is already established in the motivating paper \cite{AHLLSSV}. The point
here is the explicit separation inside a uniform ordinal--finite grid
family, together with identification of exactly which finite data survive.

\subsection{An infinite family of separations}
For $n\geq2$, let $V_n$ consist of one minimum below $n-1$ incomparable
points, and let $\Lambda_n=V_n^{\op}$. The filter polynomials are
\[
 U_{V_n}(t)=(1+t)^{n-1}+t^n,
 \qquad U_{\Lambda_n}(t)=1+t(1+t)^{n-1}.
\]
Indeed, a filter in $V_n$ either omits the minimum and is an arbitrary
subset of the top antichain, or contains everything. In the dual, a
nonempty filter must contain the unique maximum and may contain any
subset of the other points.

For every $n\geq3$, their coefficients of $t$ differ. Their powerset
types at $\omega+1$ consequently differ, although all finite order-map
counts agree by duality. The underlying grids have the common invariant
triple
\[
 (\omega\cdot n+n,\ \omega+2,\ n),
\]
by the same proof as \cref{prop:separation}.

\section{A finite signature and a universal ordinal probe}

\subsection{Ordered fiber sizes retain the missing information}
A composition of a positive integer $n$ is an ordered tuple
$a=(a_1,\ldots,a_\ell)$ of positive integers summing to $n$.
Define
\begin{equation}\label{eq:signature}
 c_Q(a_1,\ldots,a_\ell)
 :=\#\{s\in\Mon(Q,\ell): |s^{-1}(j-1)|=a_j\text{ for all }j\}.
\end{equation}
These maps are surjective. We call this finite collection of integers the
\emph{ordered-fiber signature}. There are $2^{n-1}$ compositions of $n$.
For a chain every entry is $1$; for an antichain the entry is the
multinomial coefficient $n!/(a_1!\cdots a_\ell!)$.

The same number counts chains of ideals
\[
 \varnothing=I_0\subsetneq I_1\subsetneq\cdots\subsetneq I_\ell=Q,
 \qquad |I_j\setminus I_{j-1}|=a_j.
\]
This follows by taking cumulative sublevel sets of the map. Thus the
signature is also a collection of flag counts in the finite ideal lattice;
no transfinite object is needed to compute it.

For a fixed number of blocks $m$, introduce the polynomial
\[
 K_{Q,m}(x_0,\ldots,x_{m-1})
 =\sum_{s\in\Mon(Q,m)}\prod_{q\in Q}x_{s(q)}.
\]
It remembers the ordered histogram of each map, whereas
$\Omega_Q(m)=K_{Q,m}(1,\ldots,1)$ forgets that information. The main theorem
is evaluation of this polynomial at $x_j=\omega^{\gamma_j}$, using natural
sum and natural product for its polynomial operations.

The ordinary order-map counts have the expansion
\begin{equation}\label{eq:orderpoly}
 \Omega_Q(k)=\sum_{a\vDash n}c_Q(a)\binom{k}{\ell(a)}.
\end{equation}
To see this, choose the $\ell(a)$ used colors from the $k$ available colors
and then choose the compressed surjection counted by $c_Q(a)$. This proves
polynomiality directly. It also locates the loss of information: for each
length $\ell$, the polynomial retains only the sum of signature entries
of that length, not their separate ordered fiber sizes.

\begin{proposition}[Signature evaluation]\label{prop:signatureeval}
For $n>0$ and $\beta$ as in \eqref{eq:beta},
\begin{equation}\label{eq:signatureeval}
 \otype(\Mon(Q,\beta))
 =\bigoplus_{a\vDash n}\ 
   \bigoplus_{0\leq j_1<\cdots<j_\ell<m}
   \omega^{\bigoplus_{r=1}^{\ell}\Nrep{\gamma_{j_r}}{a_r}}
       \cdot c_Q(a),
\end{equation}
where $a=(a_1,\ldots,a_\ell)$ and $a\vDash n$ means that $a$ is a
composition of $n$.
\end{proposition}
\begin{proof}
Every block assignment has an ordered set of nonempty blocks
$j_1<\cdots<j_\ell$. Compressing these labels to $0,\ldots,\ell-1$
produces a surjective isotone map. Its ordered fiber sizes form a
composition $a$, and there are exactly $c_Q(a)$ such maps. Conversely,
any such surjection can be expanded into the chosen block labels. Grouping
the terms of \eqref{eq:main} by these data gives the formula.
\end{proof}

The fork has entries
\[
 c_V(3)=1,\quad c_V(2,1)=2,\quad c_V(1,2)=1,
 \quad c_V(1,1,1)=2.
\]
For the dual, the middle two entries are interchanged. The final entry in
any poset is the number of linear extensions, while two-part entries
record ideal counts of specified sizes.

\subsection{One target suffices for all ordinal targets}
\begin{theorem}[Universal probe]\label{thm:probe}
Fix $n>0$, put $B=n+1$, and define
\begin{equation}\label{eq:probe}
 \Theta_n=\sum_{j=0}^{n-1}\omega^{B^{n-1-j}}
 =\omega^{B^{n-1}}+\cdots+\omega.
\end{equation}
For $n=1$, this means $\Theta_1=\omega$. For any two $n$-element posets
$Q$ and $R$, the following conditions are equivalent:
\begin{enumerate}[label=\textup{(\roman*)},leftmargin=2.3em]
\item $\otype(\Mon(Q,\beta))=\otype(\Mon(R,\beta))$ for every ordinal $\beta$;
\item $\otype(\Mon(Q,\Theta_n))=\otype(\Mon(R,\Theta_n))$;
\item $c_Q(a)=c_R(a)$ for every composition $a$ of $n$.
\end{enumerate}
In particular, a single ordinal value below $\omega^\omega$ encodes the
entire maximal-order-type profile of an $n$-element poset over arbitrary
ordinal targets.
\end{theorem}
\begin{proof}
Condition (i) implies (ii). To prove that (ii) implies (iii), apply the
main theorem at $\Theta_n$. A block assignment has histogram
$(h_0,\ldots,h_{n-1})$, where $0\leq h_j\leq n$ and $\sum_j h_j=n$.
Its exponent is the finite integer
\begin{equation}\label{eq:basecode}
 E=\sum_{j=0}^{n-1}h_jB^{n-1-j}.
\end{equation}
Because $h_j<B$, these are precisely base-$B$ digits, and distinct
histograms give distinct exponents. The coefficient of $\omega^E$ in
the output Cantor normal form therefore counts exactly the isotone maps
with that histogram; no collisions between different histograms occur.

For a composition $a=(a_1,\ldots,a_\ell)$, read the coefficient at the
exponent with digits
$(a_1,\ldots,a_\ell,0,\ldots,0)$. Its value is $c_Q(a)$. Hence equality
of the two probe outputs forces equality of all signature entries.
Finally, (iii) implies (i) by \cref{prop:signatureeval}, for every number
of blocks and every choice of ordinal exponents.

All exponents in \eqref{eq:probe} and \eqref{eq:basecode} are finite.
Both the probe and its output maximal type are therefore below
$\omega^\omega$.
\end{proof}

\begin{corollary}[A single grid probe]\label{cor:gridprobe}
For finite posets $P,R$ of the same positive size $n$, equality
\[
 \otype(\Grid(\Theta_n,P))=\otype(\Grid(\Theta_n,R))
\]
holds if and only if
$\otype(\Grid(\alpha,P))=\otype(\Grid(\alpha,R))$ for every ordinal
$\alpha$.
\end{corollary}
\begin{proof}
The probe is infinite, so $1+\Theta_n=\Theta_n$. Apply
\cref{thm:reduction,thm:probe} to $P^{\op}$ and $R^{\op}$. Equality of
the signatures implies equality at every target $1+\alpha$, including
finite ones. The converse follows by taking $\alpha=\Theta_n$.
\end{proof}

This theorem does not assert that the signature determines the poset up
to isomorphism. Its conclusion is exactly about the ordinal profile.
It also does not make an arbitrary uncountable ordinal computable; rather,
it isolates all dependence on the finite poset in a finite collection of
integers.

\section{Algorithms and reproducible verification}

\subsection{Direct block enumeration}
Given the strict comparison relation of $Q$ and the expanded exponents of
$\beta$, enumerate all maps $s:Q\to m$, keep the isotone ones, and collect
the monomials $\omega^{E(s)}$. The procedure is exact by
\cref{thm:block}. It performs at most $m^n$ candidate-map tests, each using
at most $n^2$ comparisons, in addition to the cost of ordinal arithmetic
and expression storage. Equal resulting exponents are combined into
finite Cantor coefficients.

This is useful as a reference implementation: its logic follows the
formula directly and differs substantially from the ideal-chain method.
It is not intended to be fast on arbitrary large finite posets.

\subsection{Dynamic programming on ideals}
Let $\Id(Q)$ be the finite lattice of ideals. A block assignment corresponds
to a weak chain
\[
 \varnothing=I_0\subseteq I_1\subseteq\cdots\subseteq I_m=Q,
 \qquad s^{-1}(j)=I_{j+1}\setminus I_j.
\]
Empty blocks are allowed and correspond to equal consecutive ideals.
Initialize $D_0(\varnothing)=1$ and all other states to $0$. For the $j$th
block exponent use
\begin{equation}\label{eq:DP}
 D_{j+1}(J)=
 \bigoplus_{\substack{I\in\Id(Q)\\I\subseteq J}}
 D_j(I)\otimes\omega^{\Nrep{\gamma_j}{|J\setminus I|}}.
\end{equation}
The result is $D_m(Q)$.

\begin{proposition}\label{prop:DP}
The recurrence \eqref{eq:DP} returns
$\otype(\Mon(Q,\beta))$.
\end{proposition}
\begin{proof}
After $j$ blocks, a state $I$ records the elements already assigned to
those blocks. The ideal condition is exactly what isotonicity requires
of a cumulative sublevel set. Extending $I$ to $J$ assigns the new block
to $J\setminus I$ and adds
$\Nrep{\gamma_j}{|J\setminus I|}$ to the exponent. Natural multiplication
of monomials performs that addition. Summing all predecessor states counts
every ideal chain once, hence every block assignment once. At the final
state all elements have been assigned.
\end{proof}

There are at most $2^n$ ideal states and at most $3^n$ ordered pairs
$I\subseteq J$ of subsets. Consequently the recurrence uses at most
$m3^n$ state transitions. This bound does \emph{not} include the cost or
size growth of symbolic ordinal expressions. The supplied implementation
precomputes the list of comparable ideal pairs by a simple quadratic scan
of the ideal list, so its preprocessing may take $O(4^n)$ subset tests;
it does not claim an overall $O(m3^n)$ bit-complexity bound.

For repeated evaluations on a fixed $Q$, the ordered-fiber signature is
another reusable representation. Formula \eqref{eq:signatureeval} then
separates the finite-poset computation from the ordinal substitution.

\subsection{Executable representation and limits}
The archive includes standard-library Python code. Its ordinal datatype
uses canonical hereditary Cantor normal forms and thus represents ordinals
below $\varepsilon_0$. A finite ordinal is a nonnegative JSON integer;
a general represented ordinal has the form
\begin{lstlisting}
{"cnf": [[exponent, positive_coefficient], ...]}
\end{lstlisting}
with recursively encoded, strictly decreasing exponents. For example,
$\omega+1$ is
\begin{lstlisting}
{"cnf": [[1, 1], [0, 1]]}
\end{lstlisting}
and $\omega^{\omega+1}$ is
\begin{lstlisting}
{"cnf": [[{"cnf": [[1, 1], [0, 1]]}, 1]]}
\end{lstlisting}
The theorem applies to more ordinals than this finite notation language.
The software is therefore an exact implementation on a specified domain,
not a universal notation system for all ordinals.

The main module validates the finite relation, closes a supplied acyclic
edge relation transitively, rejects malformed Cantor forms and negative
coefficients, and imposes explicit size limits on exponential calculations.
In grid mode it takes $P$ and $\alpha$ as input and performs the dualization
and $1+\alpha$ conversion internally. Users should not dualize the input
again.

\subsection{Checks actually performed}
The recorded verification run used Python 3.13.5 on 19 September 2026.
It enumerated all naturally labeled transitive posets on at most five
vertices: comparisons are allowed only from a smaller label to a larger
label. The counts by size $0,1,\ldots,5$ were
\[
 1,\ 1,\ 2,\ 7,\ 40,\ 357,
\]
for a total of $408$. These are not counts of isomorphism classes or of
all arbitrary labelings. Every poset of the relevant size has at least one
such topological labeling.

\begin{center}
\begin{tabular}{@{}lr@{}}
\toprule
Check category & Recorded count\\
\midrule
Finite map spaces, targets of sizes $0$ through $4$ & 2,040\\
Individual finite-map ranks recomputed from predecessors & 130,736\\
Symbolic direct-formula versus ideal-DP comparisons & 2,040\\
Finite Hoare quotient cases & 153\\
Finite subsets enumerated in those cases & 11,483\\
Distinct generated-downset classes in those cases & 2,043\\
Random ordinal-arithmetic triples, fixed seed & 400\\
Universal-probe decodings & 408\\
Recovered signature coefficients checked & 6,066\\
Signature-evaluation versus ideal-DP comparisons & 2,040\\
Failed assertions & 0\\
\bottomrule
\end{tabular}
\end{center}

For the rank tests, the code explicitly enumerated all predecessor maps
and computed their finite ranks recursively; it did not merely compare
two copies of the stated rank formula. For the Hoare tests, it enumerated
all subsets of small finite grids, formed their generated downsets, and
compared inclusion with the endpoint-map order.

The symbolic test targets included $\omega$, $\omega+1$, $\omega\cdot2$,
$\omega^2+\omega+1$, and a target with exponent $\omega+1$. The additional
profile tests used different targets, including
$\omega^2\cdot2+\omega+1$. Arithmetic checks tested associativity,
commutativity where appropriate, distributivity, monotonicity, and encoding
round trips. The JSON logs state the limits of these checks explicitly.

No finite test establishes the arbitrary-ordinal monomial theorem, the
existence of maximal linear extensions, or novelty in the literature.
Those matters must be judged from the mathematical proofs and source
comparison, not from a count of passing assertions.

\section{What the attack resolves, and what remains}

\subsection{The mathematical conclusion}
The chosen ordinal--finite family admits exact finite expressions for
both maximal order type and height. The finite factor enters the former
through its ordered-fiber signature, while the latter depends only on the
number of elements. The universal-probe theorem gives a precise
sufficiency-and-necessity statement: the signature is exactly the data
needed to recover all the maximal-order-type calculations in this family.

Three points make this more than a collection of examples. First, the
monomial proof handles mixed and arbitrary ordinal exponents rather than
only powers with finite exponents. Second, the block proof constructs a
compatible order of the fibers, so no unjustified replacement of ordinary
sum by natural sum occurs. Third, the signature theorem identifies the
finite information that the ordinary order polynomial discards.

\subsection{Boundaries of the result}
The motivating paper also considers ordinal width and broader closure
constructions \cite{AHLLSSV}. Neither a full width formula for our map
spaces nor a closure theorem under arbitrary iterations of $\Pf$ is
proved here. Computing $\otype(\Grid(\alpha,P))$ does not by itself give
$\otype(\Pf(\Grid(\alpha,P)))$: the second application need not be an
ordinal--finite grid.

Similarly, the finite signature is intrinsic to the finite factor, not an
intrinsic replacement for all weakened invariants on arbitrary wqos. An
attempt to promote it to a general invariant would need an additional
representation theorem beyond what is established here.

\subsection{Concrete next problems}
A first sharply delimited extension is to determine
$\width(\Mon(Q,\beta))$ from $Q$ and the Cantor blocks of $\beta$.
Even where maximal order type is unchanged by a monotonicity constraint,
width may be sensitive to it. The proof of the present main theorem does
not compute antichain-tree ranks.

A second problem is to identify further wpos whose finitely generated
downsets have a finite-coordinate representation with tractable fibers.
The successful ingredients here are a finite block index poset, explicit
pure-block maximal types, and an exponent weight that reverses the
pointwise order of block indices. Without an analogue of that last
compatibility, a disjoint-union upper bound need not be attained.

A third problem is to classify finite posets with the same ordered-fiber
signature, or to find smaller separating probe ordinals under additional
structural assumptions. The universal probe is designed for collision-free
encoding, not for minimality. Its base-$n+1$ exponents can be very large
finite integers; this is harmless mathematically but may be unnecessary
for particular families.

\subsection{Literature and priority assessment}
The source check on the manuscript date located the motivating extension
question and the standard maximal-order-type ingredients. Targeted searches
for maximal order types of isotone maps, finite-poset ordinal maps, and
finitary ordinal--finite grids did not establish an earlier occurrence of
the exact package of formulas proved here. Such a search is not a proof of
novelty. This manuscript supplies a detailed solution to the selected calculation
problem, not the broader extension program. Priority remains unconfirmed.

\appendix
\section{A worked three-block calculation}\label{app:worked}
Let $Q=V$ be the fork as a \emph{domain of isotone maps}, not as the finite
factor of a grid, and take
\[
 \beta=\omega^2+\omega+1.
\]
The three block exponents are $2,1,0$. If the minimum $a$ has block value
$0$, each of $b,c$ independently has one of the three block values. Their
weight contributions are $2,1,0$ and the contribution from $a$ is $2$.
The nine assignments therefore have total weights
\[
 6\ (1\text{ time}),\quad
 5\ (2),\quad 4\ (3),\quad 3\ (2),\quad 2\ (1).
\]
If $a$ has block value $1$, both upper elements have block value $1$ or
$2$, giving total weights $3$ once, $2$ twice, and $1$ once. If $a$ has
block value $2$, there is a single assignment of total weight $0$.
Collecting all $14$ assignments gives
\begin{equation}\label{eq:workedV}
 \otype(\Mon(V,\omega^2+\omega+1))
 =\omega^6+\omega^5\cdot2+\omega^4\cdot3+
   \omega^3\cdot3+\omega^2\cdot3+\omega+1.
\end{equation}
The height is independent of the fork relation. The target is a successor
of $\omega^2+\omega$, so
\[
 \height(\Mon(V,\omega^2+\omega+1))
 =\omega^2\cdot3+\omega\cdot3+1.
\]
For the corresponding grid with finite factor $\Lambda=V^{\op}$, the
same formulas apply, since the endpoint representation dualizes the finite
factor and $1+\beta=\beta$.

For the universal probe at size three,
$\Theta_3=\omega^{16}+\omega^4+\omega$, the output for $V$ is
\begin{align*}
 \otype(\Mon(V,\Theta_3))={}&
 \omega^{48}+\omega^{36}\cdot2+\omega^{33}\cdot2+
 \omega^{24}+\omega^{21}\cdot2+\omega^{18}\\
 &+\omega^{12}+\omega^9\cdot2+\omega^6+\omega^3.
\end{align*}
For instance, the exponent with base-$4$ digits $(2,1,0)$ is
$2\cdot16+1\cdot4=36$, and its coefficient $2$ is $c_V(2,1)$.
The digits $(1,2,0)$ give exponent $24$ and coefficient $1$, recovering
$c_V(1,2)$. The digits $(1,1,1)$ give exponent $21$ and coefficient $2$,
the number of linear extensions of $V$.

\section{Proof dependencies and audit points}

The proof dependencies are short enough to state explicitly. The
endpoint representation uses only finite maxima and downset inclusion.
The monomial theorem uses the standard Cartesian-product upper bound and
an explicit attaining linear extension. The general block theorem uses
the monomial theorem, the finite disjoint-union identity, and the
reverse-weight compatibility \eqref{eq:reverseweight}. The rank theorem
uses none of these maximal-order-type calculations. The height theorem
uses the rank theorem. The signature and probe theorems use the block
formula and elementary finite counting.

For an independent audit, the following are the points where a seemingly
plausible shortcut could otherwise introduce an error.
\begin{enumerate}[leftmargin=2em,itemsep=5pt]
\item \textbf{Product convention.} All wpo products in the formulas are
Cartesian. The construction inside the monomial proof concatenates
explicit extensions; it invokes no general maximal-order-type formula
for arbitrary lexicographic products.
\item \textbf{Mixed exponents.} The induction is on the number of Cantor
terms of $\delta$. Ordinary multiplication $\delta n$ is used only in
the one-term case, where it equals $\Nrep{\delta}{n}$.
\item \textbf{Fiber compatibility.} In the block theorem, earlier
comparable block assignments have larger or equal exponent weight. This
is what permits a decreasing-exponent listing that is also a linear
extension of all cross-fiber comparisons.
\item \textbf{Rank lower bound.} The coordinate that is decreased is
minimal in its equal-value fiber. Without that choice, a predecessor
with the same value could violate isotonicity after the decrease.
\item \textbf{Finite versus infinite grids.} The target is $1+\alpha$,
not $\alpha+1$, and simplifies to $\alpha$ only when $\alpha$ is infinite.
The empty generated downset is retained.
\item \textbf{Dual orientation.} The coordinate maps are antitone on
$P$, hence isotone on $P^{\op}$. In the $\omega+1$ calculation the
coefficients count filters of $P$.
\item \textbf{Probe decoding.} Histogram digits are at most $n$, strictly
less than the base $n+1$. The uniqueness argument would fail without a
collision-free choice of weights.
\end{enumerate}

\section{Archive contents and commands}
The source archive contains this \LaTeX{} file and its compiled PDF,
\texttt{README.md}, a research-status ledger, source locators, a
\texttt{Makefile}, executable Python modules, example JSON input, and the
recorded verification data. No third-party font files or copies of the
cited papers are included.

From the archive's top-level directory, run
\begin{lstlisting}
python code/ordinal_maps.py --demo
python code/ordinal_maps.py data/fork_grid_input.json
python code/verify.py
python code/profiles.py
\end{lstlisting}
The two verification scripts write their machine-readable results to
\texttt{data/}. To rebuild the paper with a suitable TeX installation,
run \texttt{make pdf}, or invoke \texttt{pdflatex} on
\texttt{ordinal\_order\_maps.tex} three times. The document uses standard
TeX packages, including \texttt{newpx}, \texttt{tcolorbox}, and
\texttt{tikz}.

\begin{thebibliography}{9}
\bibitem{AHLLSSV}
Sergio Abriola, Simon Halfon, Aliaume Lopez, Sylvain Schmitz,
Philippe Schnoebelen, and Isa Vialard.
\emph{Measuring well-quasi-ordered finitary powersets}.
arXiv:2312.14587v2, 2024.
\url{https://arxiv.org/abs/2312.14587v2}.
The motivating open direction appears in Section~6, under
``Elementary wqos.''

\bibitem{deJonghParikh}
D.~H.~J. de Jongh and Rohit Parikh.
Well-partial orderings and hierarchies.
\emph{Indagationes Mathematicae} \textbf{39}(3) (1977), 195--207.
\href{https://doi.org/10.1016/1385-7258(77)90067-1}
{doi:10.1016/1385-7258(77)90067-1}.

\bibitem{DSS}
Mirna D\v{z}amonja, Sylvain Schmitz, and Philippe Schnoebelen.
On ordinal invariants in well quasi orders and finite antichain orders.
In \emph{Well Quasi-Orders in Computation, Logic, Language and Reasoning},
Trends in Logic \textbf{53}, Springer, 2020, pp.~29--54.
Consulted preprint: arXiv:1711.00428v3, 2024.
\url{https://arxiv.org/abs/1711.00428v3}.
The standard ingredients used here are Theorem~3.4 and Lemmas~4.2,
4.5, and~4.6; no lexicographic-product identity is imported.
\end{thebibliography}
\end{document}
